% Generated by specs/tools/tested-environment.sh on 2026-10-08 22:35 UTC.
% Do not hand-edit; rerun the script on the machine that ran the tests.
\section{Tested Environment}
\label{sec:environment}

Every result in this document, and every test referenced by it, was produced on the
single machine described below, except that continuous integration also runs the Go
tests on macOS and Linux and the Chrome end-to-end suite on Linux, on GitHub's hosted
runners.  The Linux distribution tests ran on this machine too, in Docker containers
that share one Linux virtual machine (listed below).  Nothing else has been run on any
other machine, and nothing at all on the browsers this lists as not installed.  Section~\ref{sec:untested} states plainly
what that leaves unchecked.

\subsection{Hardware}

\begin{center}
\begin{tabular}{@{}ll@{}}
\toprule
Model & MacBook Pro \\
Chip & Apple M2 Max \\
Cores & 12 \\
Memory & 64~GB \\
Architecture & \code{arm64} \\
\bottomrule
\end{tabular}
\end{center}

\subsection{Operating system and browsers}

\begin{center}
\begin{tabular}{@{}>{\raggedright\arraybackslash}p{0.30\linewidth}>{\raggedright\arraybackslash}p{0.62\linewidth}@{}}
\toprule
macOS & 27.0.1 (build 26A434) \\
Safari & 27.0.1 \\
\code{safaridriver} & Included with Safari 27.0.1 (22625.1.29.11.28) \\
Google Chrome & Google Chrome 154.0.8037.98  \\
Firefox & not installed; not tested (\fq{DST-4} note; Section~\ref{sec:untested}) \\
\bottomrule
\end{tabular}
\end{center}

\subsection{Linux distributions}

The images of the last \code{make linux} (\code{e2e/linux}), with the release each
reports in \code{/etc/os-release}.  Docker 29.8.2 (engine 29.5.2), in the virtual machine of
colima version 0.10.3; every container runs on its kernel, 6.8.0-117-generic, Ubuntu 24.04.4 LTS, aarch64.  An
\code{amd64} image runs under emulation.

\begin{center}
\begin{tabular}{@{}ll@{}}
\toprule
Image & Distribution, platform \\
\midrule
\code{fsb-smoke:debian} & Debian GNU/Linux 13 (trixie), \code{linux/arm64} \\
\code{fsb-smoke:ubuntu} & Ubuntu 24.04.5 LTS, \code{linux/arm64} \\
\code{fsb-smoke:fedora} & Fedora Linux 44 (Container Image), \code{linux/arm64} \\
\code{fsb-smoke:alpine} & Alpine Linux v3.24, \code{linux/arm64} \\
\code{fsb-smoke:arch} & Arch Linux, \code{linux/amd64} \\
\code{fsb-unit:debian} & Debian GNU/Linux 13 (trixie), \code{linux/arm64} \\
\code{fsb-unit:alpine} & Alpine Linux v3.24, \code{linux/arm64} \\
\code{fsb-unit:fedora} & Fedora Linux 44 (Container Image), \code{linux/arm64} \\
\bottomrule
\end{tabular}
\end{center}

\subsection{Go toolchain and modules}

\begin{center}
\begin{tabular}{@{}>{\raggedright\arraybackslash}p{0.30\linewidth}>{\raggedright\arraybackslash}p{0.62\linewidth}@{}}
\toprule
Toolchain & \code{go version go1.27.1 darwin/arm64} \\
\code{GOOS}/\code{GOARCH} & \code{darwin}/\code{arm64} \\
\code{CGO\_ENABLED} & 1 (the program uses no cgo; \code{clang} is listed below only because the toolchain can call it) \\
\bottomrule
\end{tabular}
\end{center}

Resolved module versions (\code{go list -m all}), exactly as recorded in \code{go.sum}:

\begin{center}
\begin{tabular}{@{}ll@{}}
\toprule
Module & Version \\
\midrule
\code{golang.org/x/mod} & \code{v0.41.0} \\
\code{golang.org/x/sync} & \code{v0.23.0} \\
\code{golang.org/x/sys} & \code{v0.48.0} \\
\code{golang.org/x/text} & \code{v0.42.0} \\
\code{golang.org/x/tools} & \code{v0.49.0} \\
\bottomrule
\end{tabular}
\end{center}

\subsection{Front end: Node and npm packages}

\begin{center}
\begin{tabular}{@{}ll@{}}
\toprule
Node.js & \code{v26.11.0} \\
npm & \code{11.20.0} \\
\bottomrule
\end{tabular}
\end{center}

Installed package versions in \code{web/} (\code{npm ls --depth=0}), which match the
versions pinned in \code{web/package-lock.json}:

\begin{center}
\begin{tabular}{@{}ll@{}}
\toprule
Package & Version \\
\midrule
\code{@sveltejs/vite-plugin-svelte} & \code{7.3.1} \\
\code{@types/node} & \code{24.19.1} \\
\code{dompurify} & \code{3.4.16} \\
\code{highlight.js} & \code{11.12.0} \\
\code{marked} & \code{18.0.14} \\
\code{svelte} & \code{5.57.1} \\
\code{svelte-check} & \code{4.7.6} \\
\code{typescript} & \code{6.0.3} \\
\code{vite} & \code{8.3.1} \\
\bottomrule
\end{tabular}
\end{center}

\subsection{Build and documentation tools}

\begin{center}
\begin{tabular}{@{}>{\raggedright\arraybackslash}p{0.30\linewidth}>{\raggedright\arraybackslash}p{0.62\linewidth}@{}}
\toprule
\code{make} & GNU Make 3.81 \\
\code{git} & git version 2.56.0 \\
\code{clang} (Xcode Command Line Tools) & Apple clang version 21.0.0 (clang-2100.3.34.2) \\
\code{mandoc} & bundled with macOS 27.0.1 as /usr/bin/mandoc (BSD mandoc; the tool exposes no version flag) \\
\LaTeX{} & pdfTeX 3.141592653-2.6-1.40.29 (TeX Live 2026) \\
\code{latexmk} & Latexmk, John Collins, 9 March 2026. Version 4.88 \\
\bottomrule
\end{tabular}
\end{center}

The \LaTeX{} packages loaded by the shared preamble (\code{specs/common/preamble.tex})
are: \code{fontenc}, \code{inputenc}, \code{lmodern}, \code{geometry},
\code{microtype}, \code{parskip}, \code{amsmath}, \code{amssymb}, \code{amsthm},
\code{mathtools}, \code{booktabs}, \code{longtable}, \code{array}, \code{enumitem},
\code{needspace}, \code{xcolor}, \code{listings}, \code{fancyhdr}, \code{hyperref},
\code{bookmark}, \code{xurl} and \code{tikz}, all from the TeX Live distribution
named above.

\subsection{Regenerating this section}

This section is generated, not written by hand.  On the machine that ran the tests:

\begin{lstlisting}[language=bash]
specs/tools/tested-environment.sh > specs/design/sections/appendix-environment.tex
make -C specs design
\end{lstlisting}
